\section{Sense and Nonsense}

\begin{quotex}
We are men of the middle ages, not only because that is our destiny, the fatality of history, but also because we will it. You, you are still men of modern times, because you refuse to choose. \flright{\textsc{Nicolas Berdyaev}}

The power of a new Middle Ages is needed. \flright{\textsc{Julius Evola}, \emph{Pagan Imperialism}}

Samsara is Nirvana \flright{\textsc{Nagarjuna}}

The sensible is the symbol of the divine \flright{\textsc{St Bonaventure}}

\end{quotex}
The attack against the power of a new Middle Ages is constant, sometimes even coming from unexpected sources. Recently a bizarre argument has been circulating around the Internet. It claims that the starting point for the metaphysics of Thomas Aquinas, to wit, that knowledge begins with the senses, ``leads to" the modern world. When pressed about the meaning of ``leads to", the author has no answer. Either there is a logical or empirical necessity coercing the conclusion from the premises, or there is not.

By that logic, for example, Euclidean geometry ``leads to" Riemannian geometry. But the latter actually denies one of Euclid's postulates, so there is no logical relationship. Similarly, if the road ``leads to" Rome, it also leads to the pit when you walk off the edge. That then turns out to be the crux of the argument: the fellow does not want to follow the road to Rome. Hence, he presents this elaborate and prolix ruse to justify his own apostasy. Never mind that the Nordic-Roman Medieval Tradition encompasses much more than Thomas Aquinas. Even if he succeeds in pulling out that brick, like some giant Jenga game, the structure remains intact.

\paragraph{Being and Becoming}
In the very first paragraph of the first chapter of \emph{Revolt against the Modern World}, \textbf{Julius Evola} writes:

\begin{quotex}
In order to understand both the spirit of Tradition and its antithesis, modern civilization, it is necessary to begin with the fundamental doctrine of the two natures. According to this doctrine there is a physical order of things and a metaphysical one; there is a mortal nature and an immortal one; there is the superior realm of ``being" and the inferior realm of ``becoming." Generally speaking, there is a visible and tangible dimension and, prior to and beyond it, an invisible and intangible dimension that is the support, the source, and true life of the former. 

\end{quotex}
The alleged argument asserts Aquinas, by beginning with the senses, creates somehow a ``two story" philosophy. From that, the modern world can dispense with the second story of Being, focusing solely on the lower story of Becoming.

Obviously, that can happen only through the forgetfulness of Being.  The modern mind forgets about being, mistaking the changing world of appearances for Being itself. Opposed to that, Thomism tries to ``lead to" the awakening of Being by beginning with sense experience. Once understood, the world of becoming is intuited as a theophany of the world of Being. If not understood, then the worldview of naturalism arises, which is just the denial of the world of Being. \textbf{Julius Evola}, in \emph{Pagan Imperialism} identifies this ignorance:

\begin{quotex}
Far from exhausting itself in naturalism—as today only the ignorance or the tendentious falsification of some people are able to present it—beyond knowing the ideals of manly overcoming and of absolute liberation, in the pagan conception, the world was a living body, suffused with secret, divine, and demonic forces, with meanings and with symbols, as illustrated by that saying of \textbf{Olympiodorus}: ``the sensible expression of the invisible". 

\end{quotex}
Olympiodorus was the last of the pagan philosophers at Alexandria, but his statement applies likewise to the Christian philosophers who followed him. The modern mind displays ``ignorance" of knowledge, but the ``tendentious falsification" of the author's screed is much more sinister.

\textbf{Johann Fichte} understood this when he wrote about the ``Divine Idea of the World [that] lies at the bottom of Appearance"  and ``Sensuous appearance is the vesture, the embodiment of the Divine Idea." Goethe calls this the ``Open Secret", because it is open to all, yet seen by none. You need to choose whether or not to see this.

\paragraph{Hierarchy and Docetism}
\begin{quotex}
The authors of the Middle Ages could not imagine … the Universal Church without a pope. Because if the world is governed hierarchically, Christianity or the \emph{Sanctum Imperium} cannot be otherwise. Hierarchy is a pyramid which exists only when it is complete. \flright{\textsc{Valentin Tomberg}, \emph{Meditations on the Tarot}}

\end{quotex}
The opposite of Naturalism is Docetism. The former denies the world of Being, while the latter denies the world of Becoming. Docetism wants to inhabit some abstract spiritual reality, while denying its sensible expression in the physical world. Since the spiritual world is a Cosmic Hierarchy\footnote{\url{https://www.meditationsonthetarot.com/the-cosmic-hierarchy}}, it has a peak. Hence, to avoid Docetism, the world of Becoming, which is the symbol of the Divine, needs a peak. Hence, the denial of Thomism leads the author of that screed to the denial of the Pope. As we said, he is an apostate seeking, apparently, justification.

The vestiges of liberalism remain even in those who decry it the most. The following claim occurs so often these, it could be mistaken for a straw man. But it goes something like this: ``I have done a deep study of the Church Fathers and reached the conclusion that the Orthodox Church is most faithful to them." What could better exemplify the Protestant spirit? Or John Stuart Mill's ``marketplace of ideas", in which one evaluates competing worldviews and chooses one ``based on the evidence".

But that is to remain at a lower level of knowledge. Doxa, or opinion, is our knowledge of sense phenomena. Dianoia is rational discursive knowledge; that is, the ``weighing of evidence", disputation, and so on. But the highest level of knowledge is episteme, that is, the direct, unmediated intuition of Being. That is what the Fathers teach.

So when we arrive at that level of knowledge, we ``see" the truth of the cosmic order. Our ``choice" in that case is to submit to it. Tomberg explains:

\begin{quotex}
The vow of obedience is the practice of silencing personal desires, emotions and imagination in the face of reason and conscience; it is the primacy of the ideal [Being] as opposed to the apparent [Becoming], the nation as opposed to the personal, humanity as opposed to the nation, and God as opposed to humanity. It is the life of cosmic and human hierarchical ordering; it is the meaning and justification of the fact that there are Seraphim, Cherubim, Thrones; Dominions, Virtues, Powers; Principalities, Archangels, Angels; Priests, Knights and Commoners. Obedience is order: it is international law; it is the state; it is the Church; it is universal peace. True obedience is the very opposite of tyranny and slavery, since its root is the love which issues from faith and confidence. That which is above serves that which is below and that which is below obeys that which is above. Obedience is the practical conclusion to that which one recognises as the existence of something higher than oneself. Whosoever recognises God, obeys. 

\end{quotex}
Tomberg is describing the virtue of Piety — one of the Roman virtues as well as a gift of the Holy Spirit. Piety means the submission to whatever is higher, relinquishing the merely personal element. This sort of obedience is based not on ignorance, but on knowledge of Being; it is not coerced by an outside force, but is freely embraced.

This is the foundation of a Traditional society. How could such a society function if everyone starts ``examining the evidence" and coming to different conclusions? You could not rely on or trust anyone. It is the replacement of the knowledge of the Sacred with arbitrary personal opinion.



\flrightit{Posted on 2016-12-14 by Cologero }

\begin{center}* * *\end{center}

\begin{footnotesize}\begin{sffamily}



\texttt{O on 2016-12-15 at 11:20 said: }

But as you have pointed out elsewhere, the `order' in the West today is not a sacred Catholic order, but rather one dominated by the forces of 'the Antichrist'; and the worst is yet to come, according to a detached analysis of global trends as well as the `forecasts' of traditional prophetic data dealing with the end-times.

Especially in the case of a westerner born outside a Catholic framework, whose most recent ancestors were Protestants, for example, I cannot see why choosing the Eastern Orthodox tradition based on such an evaluation is somehow undesirable. It is inevitable for many westerners today to have to choose between traditions. The `episteme', direct knowledge that you mention has, at least in the case of most people, as its prerequisite condition that one already practices a traditional path aimed at such knowledge. Orthodox Hesychasm is a very clear example.


\hfill

\texttt{Paul on 2016-12-15 at 12:18 said: }

O — To the extent that bleak forecasts come true, they may separate out the faithful from the insincere. There is a definite rise, especially among the younger generation, of interest in and attendance at the Tridentine Latin Rite mass, as well as a return to the pre-Vatican II Catholic tradition in general. The current Pope's pointed question to the young regarding this trend, ``Why so rigid?" is being challenged by a movement which has been building steadily, including a current increase in priestly ordinations for celebrating the traditional rite. Some time ago, Pope Benedict predicted that the church would continue to empty, but that would not necessarily be a bad thing, since the true remnants would hold fast to preserve it for the future, as has happened in the past. Indeed, in another development, there has been increased discussion about Rod Dreher's ``Benedict Option," this in reference rather to Saint Benedict and his efforts to re-establish a spiritual base of human life in the aftermath of the collapsing Roman Empire. His monastic communities became the remnants that held fast for centuries. Today this has inspired Dreher who has been on the road touting ideas for similar community formations…


\hfill

\texttt{Eric on 2016-12-17 at 08:37 said: }

The local Roman Catholic mass: guitars, Protestant ``worship" songs, priest facing the people, modernist architecture, ``altar girls", communion in the hand, ``social justice" homilies, etc.

The local Eastern Orthodox Divine Liturgy: male choir singing ancient chant, priest facing the altar, Romanesque architecture and traditional iconography, reverence for the chalice, homilies focused on repentance, etc.

Where would you take your children?


\hfill

\texttt{Cologero on 2016-12-17 at 09:32 said: }

I understand, Eric. In a state of emergency, a man needs to make the best choice for his family. My only recommendation (and plea): accept it as a disorder and cool it with the harsh polemics.


\hfill

\texttt{Bill on 2016-12-17 at 13:11 said: }

Eric, if I were a father, I would talk my children to the Catholic Traditional Latin high Mass, preferably one with a professional choir singing in Latin in the organ loft. The Novus Ordo Missae, i.e., the New Order of Mass, is an ecumenical rite designed to please Protestants. If you read Pope St. Pius V's Quo primum tempore, you'll know that the Traditional Latin Mass rite is the only one that priests in the Roman part of the Catholic Church should use except for some liturgical rites that were already at least 200 years old when he published that document.

\url{http://www.catholictradition.org/Encyclicals/primum.htm}

Here's the Ottaviani Intervention that Cardinals Ottaviani and Bacci signed soon after Vatican II. For me, it carries plenty of weight because Ottaviani headed the Holy Office, now called the Sacred Congregation for the Doctrine of the Faith, during that council. Feel free to read only the summary if you want to do that. Even that little part of it may help explain why I attend only the Traditional Latin Mass. In my opinion, it's time for the Catholic Church to abolish the Novus Ordo partly because it should never have been invented.

\url{http://www.ewtn.com/library/curia/reformof.htm}


\hfill

\texttt{Eric on 2016-12-18 at 12:19 said: }

Colegero- 

My friend (and teacher), where did you see any ``harsh polemics" in my post? I'm as sick over the state of the Roman Church as anyone…and to be honest I would rather be RC. I was raised Southern Baptist but I came to Christ primarily through Roman Catholic thinkers (GKC, Nicolas Goméz Dávila at the top of the list), and I am absolutely in love with the Catholic Middle Ages. All I'm doing is reporting the ``facts on the ground" in my small city in West Texas.


\hfill

\texttt{Cologero on 2016-12-18 at 12:59 said: }

Eric, if something does not make sense, try to re-read it in a different context. The particular post concerned a fellow who did resort to unnecessary polemics. So, it was more a recommendation in general and certainly not specifically to you. I understand the ``facts on the ground". Don't be sick about anything; that is neither necessary nor helpful.


\hfill

\texttt{S. Wes on 2017-03-12 at 09:41 said: }

I've been frequenting this site for years, and was taken back by this post. It seems to be the most emotional offering by Cologero. I wish I knew the author he is speaking about, perhaps that would help me clarify what he's trying to say. I thought it might be Jay Dyer, who studied Thomism for years and has spent time recently refuting it.

I'm not sure if Cologero is defending Thomas Aquinas or the legitimacy of the papacy in this post. I can only assume he is defending Aquinas. His Evola references (who rejected the RCC), and some of his previous efforts (which I can't imagine and RCC hierarchy in history approving of) along with his attacks on exoteric constructs in general, would point to this post being a defense of Aquinas. But then he concludes about a submission to a higher authority, which again seems to suggest papal authority. Which would almost undermine all of his writings on this site, as I'm sure they were not submitted to the RCC hierarchy or their official publications. I'm just not sure what exactly he's trying to say here. Hopefully he writes more on this subject, I in no way am being critical of Cologero, I deeply respect him and his thoughts.


\hfill

\texttt{Cologero on 2017-03-18 at 16:40 said: }

@Wes, the opening statement is ``The attack against the power of a new Middle Ages is constant, sometimes even coming from unexpected sources."

So what do you suppose is being ``defended"? Then there is this:

\begin{quotex}
The vow of obedience is the practice of silencing personal desires, emotions and imagination in the face of reason and conscience; it is the primacy of the ideal [Being] as opposed to the apparent [Becoming], the nation as opposed to the personal, humanity as opposed to the nation, and God as opposed to humanity. It is the life of cosmic and human hierarchical ordering

\end{quotex}
Does that not answer the question about the ``higher authority" that ``personal desires, emotions, and imagination" are to be submitted to?


\end{sffamily}\end{footnotesize}
